\chapter{\IfLanguageName{dutch}{Conclusie}{Conclusion}}%
\label{ch:conclusie}
Deze bachelorproef onderzocht hoe realtime voertuigdata via OBD-II kan worden
verzameld, verwerkt en gevisualiseerd in een webgebaseerde monitoringapplicatie.
Daarbij werd onderzocht welke voertuigparameters beschikbaar zijn en hoe deze
gegevens efficiënt van een backend naar een frontend kunnen worden doorgestuurd.
Voor dit laatste werden REST API en WebSockets met elkaar vergeleken bij
verschillende updatefrequenties.

Tijdens het onderzoek werden vier voertuigen getest met behulp van een
Vgate VLinker FS OBD-II-adapter. Hieruit bleek dat de beschikbare
OBD-II-parameters kunnen verschillen tussen voertuigen. Hoewel OBD-II een
gestandaardiseerde interface biedt, ondersteunt niet ieder voertuig dezelfde
gegevens. Voor de ontwikkeling van een algemeen monitoringdashboard is het
daarom belangrijk om rekening te houden met ontbrekende of
voertuigspecifieke parameters.

Op basis van de verzamelde voertuigdata werd een realtime monitoringapplicatie
ontwikkeld met een FastAPI-backend en een React-frontend. De backend staat in
voor het uitlezen en verwerken van de OBD-II-data, waarna de gegevens naar de
frontend worden doorgestuurd en in het dashboard worden weergegeven. Daarnaast
werd een simulatieomgeving ontwikkeld waarmee de applicatie ook zonder een
fysiek voertuig kan worden ontwikkeld en functioneel getest.

Een belangrijk onderdeel van het onderzoek was de vergelijking tussen REST API
en WebSockets. Bij REST vraagt de frontend via polling op vaste tijdstippen
nieuwe voertuigdata op. Bij WebSockets wordt een blijvende verbinding gebruikt
waarlangs opeenvolgende gegevens naar de frontend kunnen worden gestuurd.

Beide technieken werden getest met update-intervallen van 100, 250 en
1000 milliseconden. Iedere configuratie werd gedurende ongeveer 60 seconden
getest en twee keer uitgevoerd. Hierdoor werden in totaal twaalf testruns
uitgevoerd. Tijdens deze tests werden de gemiddelde latency, P95-latency,
werkelijk behaalde updatefrequentie en het netwerkverbruik gemeten.

Uit de performantiemetingen blijkt dat WebSockets bij alle onderzochte
update-intervallen een lagere gemiddelde latency behalen dan REST. Bij een
interval van 100 ms bedroeg de gemiddelde latency 0,83 ms voor REST en
0,41 ms voor WebSockets. Bij 250 ms bedroegen deze waarden respectievelijk
1,68 ms en 0,95 ms, terwijl bij 1000 ms een gemiddelde latency van 1,29 ms
voor REST en 1,06 ms voor WebSockets werd gemeten.

De verschillen in P95-latency zijn minder eenduidig. Bij 100 ms bedroeg de
P95-latency 1,82 ms voor REST en 1,63 ms voor WebSockets. Bij 250 ms was het
verschil groter, met 6,38 ms voor REST tegenover 3,79 ms voor WebSockets.
Bij 1000 ms behaalde REST daarentegen een lagere P95-latency van 2,31 ms,
tegenover 4,50 ms voor WebSockets. Hierdoor kan niet worden gesteld dat
WebSockets in iedere situatie een lagere P95-latency behalen.

Het grootste verschil tussen beide communicatiemethoden werd vastgesteld bij
het netwerkverbruik. Bij een update-interval van 100 ms gebruikte REST
gemiddeld 14,01 KB/s, terwijl WebSockets gemiddeld slechts 1,34 KB/s
gebruikten. Bij 250 ms bedroeg het netwerkverbruik respectievelijk
5,59 KB/s en 0,58 KB/s. Bij 1000 ms werd 1,41 KB/s gemeten voor REST en
0,19 KB/s voor WebSockets.

Dit betekent dat WebSockets bij de onderzochte update-intervallen ongeveer
87 tot 90\% minder netwerkverkeer veroorzaakten dan REST. Dit verschil kan
worden verklaard door het gebruikte communicatieprincipe. Bij REST wordt voor
iedere update een afzonderlijk HTTP-request en een bijbehorende response
verstuurd. WebSockets maken daarentegen gebruik van een reeds bestaande
verbinding waarover opeenvolgende berichten kunnen worden verzonden.

Ook de werkelijk behaalde updatefrequentie werd onderzocht. Bij een ingesteld
interval van 100 ms werden gemiddeld 9,93 updates per seconde bereikt met REST
en 9,82 met WebSockets. Bij 250 ms bedroegen deze waarden respectievelijk
3,96 en 3,98 updates per seconde. Bij 1000 ms werden gemiddeld 0,98 en
1,00 updates per seconde gemeten. Beide communicatiemethoden waren binnen de
gebruikte testopstelling dus in staat om de ingestelde updatefrequenties
nagenoeg te behalen.

Een hogere updatefrequentie zorgt ervoor dat nieuwe voertuiggegevens sneller
beschikbaar kunnen worden gemaakt in het dashboard, maar verhoogt tegelijkertijd
de hoeveelheid communicatie die moet worden verwerkt. Vooral bij REST neemt
het netwerkverbruik hierdoor sterk toe. WebSockets blijven ook bij een
update-interval van 100 ms relatief efficiënt wat netwerkverbruik betreft.

Op basis van de uitgevoerde vergelijking kan worden geconcludeerd dat
WebSockets binnen de context van de ontwikkelde applicatie de meest geschikte
communicatiemethode zijn voor het continu doorsturen van realtime
OBD-II-voertuigdata. WebSockets behalen bij alle onderzochte intervallen een
lagere gemiddelde latency en veroorzaken aanzienlijk minder netwerkverkeer dan
REST, terwijl een vergelijkbare werkelijke updatefrequentie wordt behouden.

Van de onderzochte update-intervallen biedt 100 ms binnen deze testopstelling
de meest realtime weergave. Met WebSockets werden hierbij gemiddeld
9,82 updates per seconde bereikt, met een gemiddelde latency van 0,41 ms en
een netwerkverbruik van 1,34 KB/s. Hierdoor kan voertuigdata frequent worden
bijgewerkt zonder dat dit binnen de onderzochte toepassing leidt tot een sterk
toegenomen netwerkverbruik.

REST blijft echter een bruikbare oplossing voor toepassingen waarbij gegevens
minder frequent moeten worden opgehaald of waarbij continue realtime updates
niet noodzakelijk zijn. De implementatie is eenvoudig en ook bij REST werden
de ingestelde updatefrequenties nagenoeg behaald. Voor een dashboard waarin
voertuiggegevens voortdurend veranderen, bieden WebSockets binnen dit onderzoek
echter een efficiëntere oplossing.

Het onderzoek kent enkele beperkingen. De voertuigtests werden uitgevoerd met
een beperkt aantal van vier voertuigen, waardoor de resultaten rond beschikbare
OBD-II-parameters niet zonder meer naar alle voertuigen kunnen worden
veralgemeend. Daarnaast werd voor de vergelijking tussen REST en WebSockets
gebruikgemaakt van één fysieke testopstelling en één voertuig. De gebruikte
computer, OBD-II-adapter, voertuigconfiguratie en lokale netwerkopstelling
kunnen invloed hebben op de absolute meetwaarden.

Daarnaast werden frontend en backend lokaal op dezelfde computer uitgevoerd.
Hierdoor zijn de gemeten communicatielatencies zeer laag en zijn invloeden van
een extern netwerk niet meegenomen. Bij een toepassing waarbij backend en
frontend via een werkelijk netwerk met elkaar communiceren, kunnen de absolute
latencywaarden en verschillen tussen beide technologieën anders uitvallen.

Bij enkele individuele latencywaarden werden zeer kleine negatieve waarden
geregistreerd. Aangezien een negatieve communicatielatency fysiek niet mogelijk
is, worden deze beschouwd als meetartefacten van de gebruikte timestampmethode
en klokresolutie. Om deze reden werd bij de uiteindelijke evaluatie vooral
gekeken naar de gemiddelde latency en de P95-latency.

Ook de simulator blijft een vereenvoudiging van werkelijk voertuiggedrag.
Hoewel deze geschikt is om de applicatie zonder fysiek voertuig te ontwikkelen
en functioneel te testen, kunnen niet alle eigenschappen en beperkingen van
echte OBD-II-communicatie worden nagebootst. Daarom werden de uiteindelijke
performantiemetingen met echte voertuigdata uitgevoerd.

Ondanks deze beperkingen toont het onderzoek aan dat het mogelijk is om
OBD-II-data via een webgebaseerde applicatie realtime te verzamelen, verwerken
en visualiseren. De ontwikkelde oplossing brengt de volledige datastroom samen:
van het uitlezen van voertuigparameters via OBD-II tot het verwerken,
doorsturen en weergeven van deze gegevens in een realtime dashboard.

Voor toekomstig onderzoek kan de testopstelling worden uitgebreid met meer
voertuigen, langere meetperiodes en verschillende hardwareconfiguraties.
Daarnaast kan de vergelijking worden uitgevoerd wanneer frontend en backend
zich op afzonderlijke systemen en netwerken bevinden. Ook andere
communicatietechnologieën, zoals Server-Sent Events of MQTT, kunnen worden
opgenomen in een verdere vergelijking. Ten slotte kan de simulatieomgeving
worden uitgebreid met meer realistische rijsituaties en onderlinge relaties
tussen verschillende voertuigparameters.

Samenvattend blijkt binnen de onderzochte testopstelling dat WebSockets de
meest geschikte oplossing vormen voor het realtime doorsturen van
OBD-II-voertuigdata. Vooral bij een update-interval van 100 milliseconden bieden
WebSockets een combinatie van frequente updates, een lage gemiddelde latency
en aanzienlijk minder netwerkverbruik dan REST.